\documentclass[11pt]{article}
\usepackage{ochamath}
\subjectcolor{2F5DA8}
\setsheet{線形代数・電気系院試補充}{部分空間の和・共通部分・直和　演習}{https://university-math-crj.pages.dev/linear-algebra/subspace-sums/}
\namelinetrue
\begin{document}
\makesheethead{問題}
自作問題。本文の例題とは数値や設定が異なります。条件・途中式・検算を示してください。
\problem[平面の交わり]
$U=\{x\mid x_1+x_2=0\},W=\{x\mid x_2+x_3=0\}\subset\mathbb R^3$。交わりと和の基底を求めよ。
\workspace{45mm}
\problem[次元公式の証明]
$\dim(U+W)=\dim U+\dim W-\dim(U\cap W)$ を基底の延長で証明せよ。
\workspace{45mm}
\newpage
\problem[直和の一意性]
$U\cap W=\{0\}$ と、$U+W$ の各ベクトルの分解 $u+w$ が一意であることの同値を証明せよ。
\workspace{45mm}
\problem[斜めの射影]
$P=\begin{pmatrix}1&3\\0&0\end{pmatrix}$ について像・核・固有値・階数を求め、直交射影か判定せよ。
\workspace{45mm}
\newpage
\problem[冪等行列の分解]
$P^2=P$ なら $V=\operatorname{Im}P\oplus\ker P$ と $\operatorname{rank}P=\operatorname{tr}P$ を証明せよ。
\workspace{45mm}
\problem[不変部分空間]
$U$ が $A$ 不変なら、$U$ の基底を先に並べた基底で $A$ がブロック上三角になる理由を説明せよ。
\workspace{45mm}
\end{document}
